% Options for packages loaded elsewhere
% Options for packages loaded elsewhere
\PassOptionsToPackage{unicode}{hyperref}
\PassOptionsToPackage{hyphens}{url}
\PassOptionsToPackage{dvipsnames,svgnames,x11names}{xcolor}
%
\documentclass[
  letterpaper,
  DIV=11,
  numbers=noendperiod]{scrartcl}
\usepackage{xcolor}
\usepackage{amsmath,amssymb}
\setcounter{secnumdepth}{5}
\usepackage{iftex}
\ifPDFTeX
  \usepackage[T1]{fontenc}
  \usepackage[utf8]{inputenc}
  \usepackage{textcomp} % provide euro and other symbols
\else % if luatex or xetex
  \usepackage{unicode-math} % this also loads fontspec
  \defaultfontfeatures{Scale=MatchLowercase}
  \defaultfontfeatures[\rmfamily]{Ligatures=TeX,Scale=1}
\fi
\usepackage{lmodern}
\ifPDFTeX\else
  % xetex/luatex font selection
\fi
% Use upquote if available, for straight quotes in verbatim environments
\IfFileExists{upquote.sty}{\usepackage{upquote}}{}
\IfFileExists{microtype.sty}{% use microtype if available
  \usepackage[]{microtype}
  \UseMicrotypeSet[protrusion]{basicmath} % disable protrusion for tt fonts
}{}
\makeatletter
\@ifundefined{KOMAClassName}{% if non-KOMA class
  \IfFileExists{parskip.sty}{%
    \usepackage{parskip}
  }{% else
    \setlength{\parindent}{0pt}
    \setlength{\parskip}{6pt plus 2pt minus 1pt}}
}{% if KOMA class
  \KOMAoptions{parskip=half}}
\makeatother
% Make \paragraph and \subparagraph free-standing
\makeatletter
\ifx\paragraph\undefined\else
  \let\oldparagraph\paragraph
  \renewcommand{\paragraph}{
    \@ifstar
      \xxxParagraphStar
      \xxxParagraphNoStar
  }
  \newcommand{\xxxParagraphStar}[1]{\oldparagraph*{#1}\mbox{}}
  \newcommand{\xxxParagraphNoStar}[1]{\oldparagraph{#1}\mbox{}}
\fi
\ifx\subparagraph\undefined\else
  \let\oldsubparagraph\subparagraph
  \renewcommand{\subparagraph}{
    \@ifstar
      \xxxSubParagraphStar
      \xxxSubParagraphNoStar
  }
  \newcommand{\xxxSubParagraphStar}[1]{\oldsubparagraph*{#1}\mbox{}}
  \newcommand{\xxxSubParagraphNoStar}[1]{\oldsubparagraph{#1}\mbox{}}
\fi
\makeatother


\usepackage{longtable,booktabs,array}
\usepackage{calc} % for calculating minipage widths
% Correct order of tables after \paragraph or \subparagraph
\usepackage{etoolbox}
\makeatletter
\patchcmd\longtable{\par}{\if@noskipsec\mbox{}\fi\par}{}{}
\makeatother
% Allow footnotes in longtable head/foot
\IfFileExists{footnotehyper.sty}{\usepackage{footnotehyper}}{\usepackage{footnote}}
\makesavenoteenv{longtable}
\usepackage{graphicx}
\makeatletter
\newsavebox\pandoc@box
\newcommand*\pandocbounded[1]{% scales image to fit in text height/width
  \sbox\pandoc@box{#1}%
  \Gscale@div\@tempa{\textheight}{\dimexpr\ht\pandoc@box+\dp\pandoc@box\relax}%
  \Gscale@div\@tempb{\linewidth}{\wd\pandoc@box}%
  \ifdim\@tempb\p@<\@tempa\p@\let\@tempa\@tempb\fi% select the smaller of both
  \ifdim\@tempa\p@<\p@\scalebox{\@tempa}{\usebox\pandoc@box}%
  \else\usebox{\pandoc@box}%
  \fi%
}
% Set default figure placement to htbp
\def\fps@figure{htbp}
\makeatother


% definitions for citeproc citations
\NewDocumentCommand\citeproctext{}{}
\NewDocumentCommand\citeproc{mm}{%
  \begingroup\def\citeproctext{#2}\cite{#1}\endgroup}
\makeatletter
 % allow citations to break across lines
 \let\@cite@ofmt\@firstofone
 % avoid brackets around text for \cite:
 \def\@biblabel#1{}
 \def\@cite#1#2{{#1\if@tempswa , #2\fi}}
\makeatother
\newlength{\cslhangindent}
\setlength{\cslhangindent}{1.5em}
\newlength{\csllabelwidth}
\setlength{\csllabelwidth}{3em}
\newenvironment{CSLReferences}[2] % #1 hanging-indent, #2 entry-spacing
 {\begin{list}{}{%
  \setlength{\itemindent}{0pt}
  \setlength{\leftmargin}{0pt}
  \setlength{\parsep}{0pt}
  % turn on hanging indent if param 1 is 1
  \ifodd #1
   \setlength{\leftmargin}{\cslhangindent}
   \setlength{\itemindent}{-1\cslhangindent}
  \fi
  % set entry spacing
  \setlength{\itemsep}{#2\baselineskip}}}
 {\end{list}}
\usepackage{calc}
\newcommand{\CSLBlock}[1]{\hfill\break\parbox[t]{\linewidth}{\strut\ignorespaces#1\strut}}
\newcommand{\CSLLeftMargin}[1]{\parbox[t]{\csllabelwidth}{\strut#1\strut}}
\newcommand{\CSLRightInline}[1]{\parbox[t]{\linewidth - \csllabelwidth}{\strut#1\strut}}
\newcommand{\CSLIndent}[1]{\hspace{\cslhangindent}#1}



\setlength{\emergencystretch}{3em} % prevent overfull lines

\providecommand{\tightlist}{%
  \setlength{\itemsep}{0pt}\setlength{\parskip}{0pt}}



 


\KOMAoption{captions}{tableheading}
\makeatletter
\@ifpackageloaded{caption}{}{\usepackage{caption}}
\AtBeginDocument{%
\ifdefined\contentsname
  \renewcommand*\contentsname{Table of contents}
\else
  \newcommand\contentsname{Table of contents}
\fi
\ifdefined\listfigurename
  \renewcommand*\listfigurename{List of Figures}
\else
  \newcommand\listfigurename{List of Figures}
\fi
\ifdefined\listtablename
  \renewcommand*\listtablename{List of Tables}
\else
  \newcommand\listtablename{List of Tables}
\fi
\ifdefined\figurename
  \renewcommand*\figurename{Figure}
\else
  \newcommand\figurename{Figure}
\fi
\ifdefined\tablename
  \renewcommand*\tablename{Table}
\else
  \newcommand\tablename{Table}
\fi
}
\@ifpackageloaded{float}{}{\usepackage{float}}
\floatstyle{ruled}
\@ifundefined{c@chapter}{\newfloat{codelisting}{h}{lop}}{\newfloat{codelisting}{h}{lop}[chapter]}
\floatname{codelisting}{Listing}
\newcommand*\listoflistings{\listof{codelisting}{List of Listings}}
\makeatother
\makeatletter
\makeatother
\makeatletter
\@ifpackageloaded{caption}{}{\usepackage{caption}}
\@ifpackageloaded{subcaption}{}{\usepackage{subcaption}}
\makeatother
\usepackage{bookmark}
\IfFileExists{xurl.sty}{\usepackage{xurl}}{} % add URL line breaks if available
\urlstyle{same}
\hypersetup{
  pdftitle={A comment on ``Cognitive causes of `like me' race and gender biases in human language production'' (Brough et al.~2024)},
  pdfauthor={ADD NAMES},
  colorlinks=true,
  linkcolor={blue},
  filecolor={Maroon},
  citecolor={Blue},
  urlcolor={Blue},
  pdfcreator={LaTeX via pandoc}}


\title{A comment on ``Cognitive causes of `like me' race and gender
biases in human language production'' (Brough et al.~2024)}
\author{ADD NAMES}
\date{2026-09-27}
\begin{document}
\maketitle


This is a \href{https://osf.io/brxtd/files/wyxd3}{template} that you can
use to report results from reproduction studies in the social sciences.
A reproduction is the repeated analysis of a published finding using the
same data. Please ensure that the study you intend to report adheres to
the following criteria.

\textbf{Reproduction is assigned to one original study that tested at
least one identical hypothesis.}

\begin{itemize}
\item[$\square$]
  Reproduction is not ``only'' a numerical reproduction of the original
  report's results but has some added feature (e.g., different software,
  code, robustness analyses).
\item[$\square$]
  No original authors have co-authored the reproduction: While
  involvement of the original authors in the form of consultancy is
  accepted, \emph{Replication Research} does not publish reproductions
  co-authored by any of the original author.
\item[$\square$]
  A dataset from the original study is analysed. If \emph{different}
  data is used, please use the template for replication studies instead:
  \url{https://osf.io/3jgxd}.
\item[$\square$]
  This is the most recent version of this document (check here:
  \url{https://osf.io/wyxd3}).
\end{itemize}

To complete this template, please provide the information underneath
each heading for your study. You need to use a single template for each
reproduction study that you want to report. The paper has four main
sections: (1) Background; (2) Methods; (3) Results; (4) Evaluation. This
template is heavily based on the data paper template by Evans (2024).
Some parts (e.g., reproduction closeness) are based on suggestions by
Brandt et al. (2014). Before publishing your report, please delete the
instructions that are written in blue.

Completed reports can be submitted to the non-peer reviewed
\href{https://zenodo.org/communities/r2/}{Zenodo Replication Research
Community} or to the peer reviewed journal
\href{https://replicationresearch.org}{Replication Research}.

\subsection*{License}\label{license}
\addcontentsline{toc}{subsection}{License}

If you are using this template (Röseler et al. 2025) or parts of this
template to structure your report, please cite it:

Röseler, L., Hein, M., \& Oppong Boakye, P. (2025). StaRT Reports.\\
https://doi.org/10.17605/OSF.IO/BRXTD

\newpage{}

\subsection*{Title}\label{title}
\addcontentsline{toc}{subsection}{Title}

\subsection*{Abstract}\label{abstract}
\addcontentsline{toc}{subsection}{Abstract}

The original study by Brough et al.~{[}-@brough2024{]} investigates how
social biases may influence automatic word retrieval mechanisms during
sentence production. Across four experiments, the authors observe
significant `like me' effects, whereby speakers preferentially mention
people who share their own race and gender before naming out‑group
members (β\,=\,0.16,\,SE\,=\,0.02, z~=\,7.7, p\,\textless\,0.001).
Psycholinguistic manipulations show that the effect does not vary across
grammatical constructions (χ²(2)\,=\,0.65, p =\,0.721).

We conducted a computational reproduction using the authors' original
raw data and R scripts, supplemented by reconstruction of some missing
scripts/steps, and performed multiple robustness checks (alternative
model specifications, alternative exclusion criteria, etc.). Full
reproduction of the original statistics was achieved for Studies\, 1--3.
The statistical results of Study\,4 were also reproduced except for a
sign error.

We conclude that the findings reported in Brough et al.~(2024) are
computationally reproducible using the authors' scripts with some minor
adjustments and additions. We also found them to be robust to our
modifications. To facilitate future repetitive research and cumulative
endeavours, we recommend that multi‑study papers provide detailed README
files.

(172 words)

\subsection*{Keywords}\label{keywords}
\addcontentsline{toc}{subsection}{Keywords}

like me biases; human language production;

\newpage{}

\section{Background}\label{background}

In ``Cognitive causes of `like me' race and gender biases in human
language production'', Brough et al.~(2024) describe how speakers tend
to first mention individuals who resemble them or their social group.
The paper claims that this `like me' effect is due to ``(\ldots)
automatized mechanisms by which humans retrieve words from memory to
produce sentences'' ((Brough et al. 2024a), p.~1706). They demonstrate
that the names of individuals who are more similar to the participants
are more accessible to the participants.

The focus of the paper is the influence of gender and race. The effect
was observed in four different studies. Testing both English and Chinese
L1 and L2 speakers, the authors claim that the `like me' effect is
independent of demographic background. In Studies 1--3, participants had
to learn the names of eight persons shown on photographs, who were
either Black or white and male or female. They then had to describe
events in which two of the persons, differing in either gender or race,
engaged in symmetrical interactions. The aim of the studies was to
measure ``the degree to which participants framed their descriptions in
ways that highlighted their own social group, by beginning their
sentences with the name of the person who matched them in gender or in
race'' ((Brough et al. 2024a), p.~1707).~

Although the overall `like me' effect was small, with mean values of
around 55\% ((Brough et al. 2024a)p.~1711), the authors conclude that
the observed social bias is important. They suggest that the effect is
independent of language, linguistic construction, and social group,
although Bough et al.~(2024: 1711) acknowledge that the first conclusion
is based on just two languages (English and Chinese). According to the
authors, ``the fact that the like me effect was invariant across
constructions suggests that it is driven by accessibility and
incrementality, rather than by evaluation'' ((Brough et al.
2024a)p.1711). They claim that ``a core conclusion from the present
studies is that the incremental and accessibility-driven language system
that allows speakers to produce language that is fluent also leads them
to produce language that is biased'' ((Brough et al. 2024b) p.~1711).
Other researchers hope that the large phenomenon of implicit biases will
be studied more in the future (Axt et al. 2025) or that
counterstereotypical AI can be used to cause a shift to no longer
represent certain group members based on stereotypical information which
is stored in the memory (Hermann, De Freitas, and Puntoni 2025).

The recent study ``Exploration of Animacy in the''Like-Me'' Effect''
(Halpern, n.d.a) aimed at replicating the original study by Brough et
al.~(2024) ``to see if the Like-Me Effect would occur on the axis of
gender, if the characters presented were animated.'' (Halpern, n.d.a).
This study failed to replicate a significant overall ``Like-Me'' effect,
but found a significant effect of Gender. Male participants were more
likely to produce a ``Like-Me'' bias than female participants (Halpern,
n.d.b). Moreover they found that men were more likely to produce a
``Like-Me'' bias in the animated condition whereas there was not a
difference between women and men in the Mixed condition to produce a
``Like-Me'' response(Halpern, n.d.c). One limitation of the replication
is, that the sample was not diverse and consisted of 32 students who
were all registered in the psychology department at Florida State
University (Halpern, n.d.b).

\subsection{Motivation for the present
study}\label{motivation-for-the-present-study}

To our knowledge the study has not been reproduced. The original study
provides open access to their data, stimuli and R scripts as well as to
their preregistrations and reviews.

Moreover, as we did have access to the preregistrations, we noticed that
the reported models in the paper for Study 2 differ considerably from
the models that were preregistered and that the preregistered models for
Study 3 differ from the models reported in the article in that they
include gender and race as fixed effects. Therefore we performed
robustness checks by computing the preregistered models using the
provided data.

In Study 4, the authors indicate that the observed `like me' gender
effects may be smaller in Chinese because of including two possibly
gender neutral chinese names ( 霍普 {[}Hope{]} and 贝斯 {]}Beth{]}).
Therefore we decided to rerun the analysis for Study 4 with a subset of
the data that excluded these two neutral names. We noticed that the main
claim made on the basis of these data relies on a GLMM in which the
`like me' effect is modelled as a numeric predictor ranging from -1 to
0. However, this variable only has three possible values that, somewhat
counterintuitively, correspond to the following:~

\begin{itemize}
\item
  -1 = highest me-likeness (i.e.~both race and gender match participant
  demographic)
\item
  0 (only race or gender matches participant demographic)
\item
  1 = lowest me-likeness (neither race nor gender match participant
  demographic
\end{itemize}

We believe that this variable would be better modelled as an ordinal
categorical variable, hence we reran the reported model with the
independent variable modelled as an ordered factor.

\subsection{Original and Reproduction Study
References}\label{original-and-reproduction-study-references}

\begin{longtable}[]{@{}lll@{}}
\toprule\noalign{}
& Reference with DOI & Study number \\
\midrule\noalign{}
\endhead
\bottomrule\noalign{}
\endlastfoot
Original & & \\
Study number & & \\
\end{longtable}

\subsection{Reproduced Study, Hypothesis, and
Findings}\label{reproduced-study-hypothesis-and-findings}

\section{Methods}\label{methods}

\subsection{Reproduction type}\label{reproduction-type}

\section{\texorpdfstring{\textbf{2. Computational
Reproducibility}}{2. Computational Reproducibility}}\label{computational-reproducibility}

The original data and code were made available on the OSF
(\url{https://osf.io/87vew/overview}). The repository includes raw and
analysis data, as well as cleaning and analysis code in R. As the
analysed paper comprises a multitude of studies, we found the total
number of subdirectories and files overwhelming. The absence of a ReadMe
file made navigating the data packages challenging.~

In the preprocessing scripts for Studies~1--3
(\texttt{Process\_Study1\_Data.R}, \texttt{Process\_Study2\_Data.R},
\texttt{Process\_Study3\_Data.R}) the authors mention that, due to a
technical issue involving an unbacked-up computer, the authors no longer
have access to all the initial raw data for these studies. However, they
provide full access to data files, which have apparently only been
hand-corrected for minor spelling mistakes (e.g.~`Rth' instead of
`Ruth'). We therefore classified these files as ``raw data
provided''.~\\

\begin{longtable}[]{@{}llll@{}}
\toprule\noalign{}
\endhead
\bottomrule\noalign{}
\endlastfoot
& \textbf{Fully} & \textbf{Partial} & \textbf{No} \\
Raw data provided & x & & \\
Cleaning code provided & x & & \\
Analysis data provided & x & & \\
Analysis code provided & x & & \\
Reproducible from raw data & & x & \\
Reproducible from analysis data & & x & \\
\end{longtable}

We were able to computationally reproduce the results of Studies~1-3
with the code provided by the authors. The only noticeable difference
was a slight change in the numbers of the 95\% bootstrapped confidence
intervals with a maximal difference of ±0.2, which we believe is due to
the way random seeds were set in the R scripts. The function set.seed(),
by default, in R only applies to the next operation that involves
randomisation. The authors only set a seed at the top of the manuscript
and not every time randomisation was intended .

Concerning Study~4, we were unable to fully reproduce two processed data
files (\texttt{study5\_chinese.csv} and \texttt{data5\_chinese.csv})
from the raw data using the corresponding preprocessing script
(\texttt{Data\ Preprocessing/Chinese\ Study/Process\_Chinese\_Data.R}).
The script returned an error because the column \texttt{TrialType} does
not exist. We noted that the raw data does include a column entitled
\texttt{trial\_type} which, however, only contains a single unique
value: ``survey\_text''. By contrast, the \texttt{Trial\_Type} variable
contained in the processed data file uploaded by the authors,
\texttt{data5\_chinese.csv}, has two levels, differentiating between
``RaceTrial'' and ``GenderTrial''. We were unable to reproduce this
variable on the basis of the raw data, but all other variables were
successfully reproduced. For Study~4, we therefore had to use the
\texttt{study5\_chinese.csv} and \texttt{data5\_chinese.csv} processed
data files provided by the authors to reproduce the statistical analyses
based on these data.~

Moreover, we were unable to reproduce the median time taken to complete
inputting the ``critical sentence'', which according to the paper is
9.5~s in the English condition and 8.3~s in the Chinese condition
((Brough et al. 2024a), p.~1713 ). It was not clear to us what exactly
was meant by ``critical sentence'' and how these could be identified in
the datasets. We also observed a minor difference in the final sample
size for women in the English condition for Study~4. The paper refers to
28 women ((Brough et al. 2024a), p.1713), but after excluding those with
slow RTs, the script \texttt{Process\_Chinese\_Data.R} contains a table
with `condition`, `ppt\_gender`, `list` and `n` in which only 27 English
women are listed.

We were able to reproduce the results of the `like me' effect reported
as part of Study~4. However, we were unable to locate the code necessary
to reproduce some of the values reported in the paper and therefore had
to write the code ourselves to check these numbers. This was the case
for size of the overall `like me' effect, reported as β = 0.18 (0.04), z
= 5.0, P \textless{} 0.001, which we successfully reproduced, as well as
the claim that ``the size of the effect did not significantly differ
depending on the language used'' ((Brough et al. 2024c), p.~1710), for
which we reproduced the same statistics as in the original study apart
from one value: we obtained a negative number for the \emph{z}-value
(\emph{z}~=~-1.299) whereas, in the paper, a positive \emph{z}-value
(\emph{z}~=~1.3) is reported ((Brough et al. 2024c), p.~1710). This sign
error is likely to be a typographical error and does not change the
claims made in the paper. As with Studies~1-3, we also observed minor
differences in the values of the 95\% bootstrapped confidence
intervals.~

For the Empathy Rating Study, the preprocessing script
(\texttt{Data\ Preprocessing\ /EmpathyStudy/process\_data.R}) filtered
participants, but did not create the processed dataset used in
subsequent analyses. Consequently, it was originally unclear to us
whether the data used in the analysis script
(\texttt{Data\ Analysis/empathy\_study\_data.csv}) matches the output of
the preprocessing steps. However, we were eventually able to
successfully recreate the dataframe used by the original authors, based
on an modified version of their preprocessing script, the information
provided in the preregistration
(\url{https://aspredicted.org/7698v.pdf}), and the raw data provided by
the authors
(\texttt{Data\ Preprocessing\ /EmpathyStudy/full\_empathy\_data.csv}).

The Agency Rating Study is based on two separate preregistered studies
and data collections (one for transitive and one for conjoint
constructions), both carried out on Prolific. We were able to
successfully reproduce the combined processed dataset from the raw
Prolific data using the preprocessing script provided
(\texttt{Data\ Preprossessing/Rating\ Study/process\_rating\_data.R}).~

We were also able to successfully reproduce all of the plots presented
in the paper (see \texttt{OSF\_Analyses\_reproduction.qmd}).

On p.~1709, the authors (Brough et al. 2024d) write:

\begin{quote}
As Fig. 2a shows, the sizes of the like me effects varied only minimally
across the grammatical constructions that we had tested (χ2(2) = 0.65, P
= 0.721). Rated agency and rated identification, by contrast, strongly
varied based upon the grammatical constructions (both χ2(2)
\textgreater{} 200, P \textless{} 0.001) and clearly dissociated from
one another.
\end{quote}

We were unable to reproduce the first chi-square value. Our
understanding of the code is that the ``Test of grammar'' section (from
line 1022 of the authors' \texttt{OSF\_Analyses.R} script) best
corresponds to this analysis. This test also returns a non-significant
\emph{p}-value (0.4886), but a chi-square value of 2.4272. We were able
to reproduce the second chi-square value in the quoted paragraph above
as 201.99 for the rated agency (transitive and conjoint sentences) and
820.96 for the rated identification (referred to in the code/data as the
Empathy Rating Study). Both correspond to \emph{p}-values that are,
indeed, \textless{} 0.001.

We successfully recreated the processed dataset necessary for the
analysis of the memory naming task (\texttt{studies1\_2\_3\_memory.csv})
based on the raw data using the script
\texttt{overal\_memory\_data\_processing.R}; however, we noted that this
script -- which is over 1000 lines long -- appears to repeat many of the
data processing steps completed in earlier data processing scripts. This
is particularly true of the removal of trials and participants according
to the exclusion criteria; the latter of which is achieved by filtering
out hard-coded participant IDs, which were presumably identified via the
visual analysis of the tables produced in the script. It also includes
some interactive elements (where the user needs to change a variable
multiple times to check various subsets of the data) and many repeated
code chunks within the script, too, making it difficult to check for
consistency. In spite of the presence of comments in the code, it is
difficult to follow the logic of the various data wrangling steps. As
delivered, this preprocessing script is not reproducible.

The main result of the memory naming experiment could be reproduced (β =
0.02 (0.005), t(642) = 4.5, p \textless{} 0.001). The mean reaction
times reported on pp.~1709--1710 could all be reproduced. The
bootstrapped confidence intervals differed slightly, but this is not
surprising as no random seed was set in the
\texttt{overall\_memory\_analysis.R} script.

For the Multi-Study analysis, which combines the data across Studies 1
through 3, the Methods section ((Brough et al. 2024d)p.~1712ff.),
describes what is being investigated in that part of the paper:~

\begin{quote}
We analysed the effect of syntax across the first three studies using a
regression of the form LikeMe \textasciitilde{} 1 + Grammatical Frame +
(1\textbar subject) + (1\textbar event).

{[}\ldots{]}

The main omnibus analyses used mixed-effects regressions of the form
LikeMe \textasciitilde{} 1 + Participant Gender × Participant Race +
(1\textbar subject) + (1 + Participant Gender × Participant
Race\textbar event). Predictors were contrast coded, and separate models
were fit to the race and gender trials. The analysis of trial type
(using data from all four production studies) had the form LikeMe
\textasciitilde{} 1 + Trial Type {[}Gender versus Race-focused{]} + (1 +
Trial Type\textbar subject) + (1 + Trial Type\textbar event)''
(p.~1713f.).
\end{quote}

These regression models can be found in the code provided by the auhors
\texttt{R\_scipts/scripts\_and\_data\_original/Data\_Analysis\_Original/OSF\_Analyses.r}
( ) and match the descriptions in the paper. We successfully reproduced
the results of these analyses.

\subsection{Original and Reproduction Study
Features}\label{original-and-reproduction-study-features}

\textbf{Table 1} \emph{Differences Between the Original and the
Reproduction Study}

\begin{longtable}[]{@{}
  >{\raggedright\arraybackslash}p{(\linewidth - 6\tabcolsep) * \real{0.2500}}
  >{\raggedright\arraybackslash}p{(\linewidth - 6\tabcolsep) * \real{0.2500}}
  >{\raggedright\arraybackslash}p{(\linewidth - 6\tabcolsep) * \real{0.2500}}
  >{\raggedright\arraybackslash}p{(\linewidth - 6\tabcolsep) * \real{0.2500}}@{}}
\toprule\noalign{}
\begin{minipage}[b]{\linewidth}\raggedright
Study Feature
\end{minipage} & \begin{minipage}[b]{\linewidth}\raggedright
Original Study
\end{minipage} & \begin{minipage}[b]{\linewidth}\raggedright
Reproduction Study
\end{minipage} & \begin{minipage}[b]{\linewidth}\raggedright
Justification and Possible Effect on Interpretability
\end{minipage} \\
\midrule\noalign{}
\endhead
\bottomrule\noalign{}
\endlastfoot
Software and versions & & & \\
Functions & & & \\
\ldots{} & \ldots{} & \ldots{} & {[}press Tab here to add new rows{]} \\
\end{longtable}

\subsection{Involvement of Original
Authors}\label{involvement-of-original-authors}

Did original authors provide:

\begin{itemize}
\item
  Code/data? → Permitted (non-authorship)~
\item
  Analysis input? → Permitted (non-authorship)~
\item
  Draft revisions? → NOT permitted (co-authorship)
\end{itemize}

\newpage{}

\subsection{Preregistration}\label{preregistration}

\begin{longtable}[]{@{}ll@{}}
\toprule\noalign{}
Study & Preregistration link \\
\midrule\noalign{}
\endhead
\bottomrule\noalign{}
\endlastfoot
Original & \\
Study number & \\
\end{longtable}

\subsection{Availability of Materials}\label{availability-of-materials}

\section{Results}\label{results}

\subsection{Original key statistical result and
interpretation}\label{original-key-statistical-result-and-interpretation}

\begin{longtable}[]{@{}
  >{\raggedright\arraybackslash}p{(\linewidth - 4\tabcolsep) * \real{0.3333}}
  >{\raggedright\arraybackslash}p{(\linewidth - 4\tabcolsep) * \real{0.3333}}
  >{\raggedright\arraybackslash}p{(\linewidth - 4\tabcolsep) * \real{0.3333}}@{}}
\toprule\noalign{}
\begin{minipage}[b]{\linewidth}\raggedright
Claim
\end{minipage} & \begin{minipage}[b]{\linewidth}\raggedright
Key results and interpretation
\end{minipage} & \begin{minipage}[b]{\linewidth}\raggedright
Reproduction result
\end{minipage} \\
\midrule\noalign{}
\endhead
\bottomrule\noalign{}
\endlastfoot
Please provide the original & & \\
authors' claim and page number & Key statistical result from the
original report & \\
Correlations between variables of interest & Table 1 & Table 2 \\
Relationship between variable x, y, and z & Figure 4 & Figure 2 \\
\end{longtable}

\subsection{Comparison of Results Between Original Study and
Reproduction}\label{comparison-of-results-between-original-study-and-reproduction}

\subsection{Further results}\label{further-results}

\subsection{Dataset and Materials
link}\label{dataset-and-materials-link}

\subsection{Object/file name}\label{objectfile-name}

\subsection{Limits to sharing}\label{limits-to-sharing}

\subsection{FAIR data/Codebook}\label{fair-datacodebook}

\begin{itemize}
\item[$\square$]
  Data deposited in certified repository (e.g., OSF, Zenodo)
\item[$\square$]
  Dataset has persistent identifier (DOI)
\item[$\square$]
  Code shared with open license (e.g., MIT, GPL)
\item[$\square$]
  Codebook includes:

  \begin{itemize}
  \tightlist
  \item[$\square$]
    Variable names Description
  \item[$\square$]
    Valid ranges
  \item[$\square$]
    Missing codes
  \end{itemize}
\item[$\square$]
  Metadata uses standard schema (e.g., Datacite, CodeMeta)
\item[$\square$]
  File formats are non-proprietary (e.g., .csv, .txt, .R)
\end{itemize}

\subsection{Deviations from the
Preregistration}\label{deviations-from-the-preregistration}

\begin{longtable}[]{@{}
  >{\raggedright\arraybackslash}p{(\linewidth - 4\tabcolsep) * \real{0.3514}}
  >{\raggedright\arraybackslash}p{(\linewidth - 4\tabcolsep) * \real{0.2973}}
  >{\raggedright\arraybackslash}p{(\linewidth - 4\tabcolsep) * \real{0.3514}}@{}}
\toprule\noalign{}
\begin{minipage}[b]{\linewidth}\raggedright
Description of Deviation
\end{minipage} & \begin{minipage}[b]{\linewidth}\raggedright
Reason for Deviation
\end{minipage} & \begin{minipage}[b]{\linewidth}\raggedright
Influence on the Results
\end{minipage} \\
\midrule\noalign{}
\endhead
\bottomrule\noalign{}
\endlastfoot
\ldots{} & \ldots{} & \ldots{} \\
& & \\
\end{longtable}

\subsection{FORRT Replication Database (FReD)
Entry}\label{forrt-replication-database-fred-entry}

Please disclose whether your results have already been entered into the
replication database. To check it, you can enter the DOI of the target
study here: \url{https://forrt.org/apps/fred_annotator.html}. If they
are, we will link this report to the entry. If they are not, we invite
you to enter them in the spreadsheet or in case of a low number of
replication effects here:
\url{https://docs.google.com/forms/d/e/1FAIpQLSeMCwdtP0TPgL55stniuyyTxnNwyC34mO4VUuLcQwYrLI89sQ/viewform}.

\section{Evaluation}\label{evaluation}

\subsection{Feedback from the original study's authors
(optional)}\label{feedback-from-the-original-studys-authors-optional}

You may reach out to the original authors. In this case, we recommend
adhering to the ``golden rule of reproductions: treat others and their
work, as you would like others to treat you and your work!'' (Berkeley
Initiative for Transparency in the Social SciencesTeam 2020). For email
templates, see the ACRE guide:
\url{https://bitss.github.io/ACRE/comunications.html\#for-reproducers-contacting-the-authors-of-the-original-study}

Original authors were contacted on {[}date{]}.\\
Their response:\\

\begin{itemize}
\item[$\square$]
  Comments appended
\item[$\square$]
  Declined to comment
\item[$\square$]
  No response after 4 weeks.
\end{itemize}

If no response, editors may independently solicit feedback.

\subsection{Limitations and Recommendations for Future
Replications}\label{limitations-and-recommendations-for-future-replications}

\subsection{Brief Evaluation}\label{brief-evaluation}

\newpage{}

Other information required for submission, not for review

\section*{Contribution Statement}\label{contribution-statement}
\addcontentsline{toc}{section}{Contribution Statement}

\section*{Ethical Statement}\label{ethical-statement}
\addcontentsline{toc}{section}{Ethical Statement}

\section*{Acknowledgements}\label{acknowledgements}
\addcontentsline{toc}{section}{Acknowledgements}

Please add any relevant acknowledgements to anyone else who supported
the project in which the data was collected, but did not work directly
on the data itself.

\section*{Conflicts of interest}\label{conflicts-of-interest}
\addcontentsline{toc}{section}{Conflicts of interest}

are defined as ``when authors, reviewers, or editors have interests that
are not fully apparent and that may influence their judgments on what is
published. They have been described as those which, when revealed later,
would make a reasonable reader feel misled or deceived'' (COPE, 2003).
Please note any potential conflict of interests, including but not
limited to direct or indirect financial gains or non-financial
implications. If you have no conflicts to declare, please state
explicitly: ``The author(s) declare no conflict of interest associated
with the publication of this manuscript.

\section*{Funding statement}\label{funding-statement}
\addcontentsline{toc}{section}{Funding statement}

If the data, or the project from which it came, required funding, please
provide clear detail of this here. For example, if funded by a research
council this would include the year of successful receipt of funding,
the name of the funding council, and the grant number.

\newpage{}

\section*{References}\label{references}
\addcontentsline{toc}{section}{References}

Please enter references (preferably in the APA style) and include a DOI
where available.

\phantomsection\label{refs}
\begin{CSLReferences}{1}{0}
\bibitem[\citeproctext]{ref-axt2025}
Axt, Jordan, Valentina Palacio Posada, Eliane Roy, and Jeffrey To. 2025.
{``Revisiting the Policy Implications of Implicit Social Cognition.''}
\emph{Social Issues and Policy Review} 19 (1): e70003.
\url{https://doi.org/10.1111/sipr.70003}.

\bibitem[\citeproctext]{ref-berkeleyinitiativefortransparencyinthesocialsciencesteam2020}
Berkeley Initiative for Transparency in the Social SciencesTeam. 2020.
\emph{Guide for Advancing Computational Reproducibility in the Social
Sciences}. \url{https://bitss.github.io/ACRE/}.

\bibitem[\citeproctext]{ref-brandt_replication_2014}
Brandt, Mark J., Hans IJzerman, Ap Dijksterhuis, Frank J. Farach, Jason
Geller, Roger Giner-Sorolla, James A. Grange, Marco Perugini, Jeffrey R.
Spies, and Anna van 't Veer. 2014. {``The {Replication} {Recipe}: {What}
Makes for a Convincing Replication?''} \emph{Journal of Experimental
Social Psychology} 50 (January): 217--24.
\url{https://doi.org/10.1016/j.jesp.2013.10.005}.

\bibitem[\citeproctext]{ref-brough2024}
Brough, Jessica, Lasana T. Harris, Shi Hui Wu, Holly P. Branigan, and
Hugh Rabagliati. 2024a. {``Cognitive Causes of {`}Like Me{'} Race and
Gender Biases in Human Language Production.''} \emph{Nature Human
Behaviour} 8 (9): 1706--15.
\url{https://doi.org/10.1038/s41562-024-01943-3}.

\bibitem[\citeproctext]{ref-brough2024a}
---------. 2024c. {``Cognitive Causes of {`}Like Me{'} Race and Gender
Biases in Human Language Production.''} \emph{Nature Human Behaviour} 8
(9): 1706--15. \url{https://doi.org/10.1038/s41562-024-01943-3}.

\bibitem[\citeproctext]{ref-brough2024b}
---------. 2024d. {``Cognitive Causes of {`}Like Me{'} Race and Gender
Biases in Human Language Production.''} \emph{Nature Human Behaviour} 8
(9): 1706--15. \url{https://doi.org/10.1038/s41562-024-01943-3}.

\bibitem[\citeproctext]{ref-brough2024e}
---------. 2024b. {``Cognitive Causes of {`}Like Me{'} Race and Gender
Biases in Human Language Production.''} \emph{Nature Human Behaviour} 8
(9): 1706--15. \url{https://doi.org/10.1038/s41562-024-01943-3}.

\bibitem[\citeproctext]{ref-evans_journal_2024}
Evans, Thomas R. 2024. {``Journal of {Open} {Psychology} {Data} - {Data}
{Paper} {Template}.''} PsyArXiv.
\url{https://doi.org/10.31234/osf.io/9ms7h}.

\bibitem[\citeproctext]{ref-halpern}
Halpern, Rebekka Jeanne Elisabeth. n.d.a. {``Exploration of Animacy in
the {"}Like-Me{"} Effect.''}
\url{https://doi.org/10.33009/FSU_f151ee87-6323-450c-a8df-c6e26ae41025}.

\bibitem[\citeproctext]{ref-halperna}
---------. n.d.b. {``Exploration of Animacy in the {"}Like-Me{"}
Effect.''}
\url{https://doi.org/10.33009/FSU_f151ee87-6323-450c-a8df-c6e26ae41025}.

\bibitem[\citeproctext]{ref-halpernb}
---------. n.d.c. {``Exploration of Animacy in the {"}Like-Me{"}
Effect.''}
\url{https://doi.org/10.33009/FSU_f151ee87-6323-450c-a8df-c6e26ae41025}.

\bibitem[\citeproctext]{ref-hermann2025}
Hermann, Erik, Julian De Freitas, and Stefano Puntoni. 2025. {``Reducing
Prejudice with Counter{-}Stereotypical AI.''} \emph{Consumer Psychology
Review} 8 (1): 75--86. \url{https://doi.org/10.1002/arcp.1102}.

\bibitem[\citeproctext]{ref-roseler_handbook_2025}
Röseler, Lukas, Lukas Wallrich, Helena Hartmann, Luisa Altegoer,
Veronica Boyce, Sarahanne M. Field, Janik Goltermann, et al. 2025.
{``Handbook for {Reproduction} and {Replication} {Studies}.''}
\url{https://forrt.org/replication-handbook/}.

\end{CSLReferences}




\end{document}
